% ACL 2023 style — use with acl.sty from https://github.com/acl-org/acl-style-files
% Compile with: pdflatex -> bibtex -> pdflatex -> pdflatex
% Or paste into Overleaf and select ACL 2023 template

\documentclass[11pt]{article}
\usepackage{amsmath}
\usepackage{amssymb}
\usepackage[hyperref]{acl}
\usepackage{times}
\usepackage{latexsym}
\usepackage{graphicx}
\usepackage{booktabs}
\usepackage{amsmath}
\usepackage{multirow}
\usepackage{url}
\usepackage{microtype}

\title{Parameter-Efficient Fine-Tuning of ModernBERT:\\
A Comparative Study of LoRA and IA\textsuperscript{3}}

\author{
  Jonna Bauer \quad Francesca Romana Blanda \quad Ting Hsu \\
  Theory and Applications of Deep Learning (IIE7721) \\
  Yonsei Uiversity \\
}

\begin{document}
\maketitle

% ─────────────────────────────────────────
\begin{abstract}
Fine-tuning large pre-trained language models for downstream tasks is
computationally expensive, requiring updates to hundreds of millions of
parameters. Parameter-efficient fine-tuning (PEFT) methods address this
by training only a small number of additional parameters while keeping
the base model frozen. In this work, we apply two prominent PEFT
methods --- Low-Rank Adaptation (LoRA) and Infused Adapter by Inhibiting
and Amplifying Inner Activations (IA\textsuperscript{3}) --- to
ModernBERT \cite{warner2024modernbert}, a recently released encoder-only
transformer that modernises the original BERT architecture. We evaluate
both methods on two text classification benchmarks: SST-2 (sentiment
analysis) and QNLI (question-answer natural language inference). Our
experiments show that LoRA achieves within 2.3\% of full fine-tuning
accuracy while using 99.6\% fewer trainable parameters, but degrades
significantly in low-data settings. We further analyse the effect of
LoRA rank and training set size on model performance, providing
practical guidance for applying PEFT methods to ModernBERT.
\end{abstract}

% ─────────────────────────────────────────
\section{Introduction}

Pre-trained language models such as BERT \cite{devlin2019bert} have
become the backbone of modern natural language processing (NLP). These
models are typically pre-trained on large corpora and then fine-tuned on
task-specific datasets. However, full fine-tuning updates all model
parameters --- hundreds of millions in the case of BERT-base --- making
it memory-intensive and expensive, especially when fine-tuning for
multiple downstream tasks simultaneously.

Parameter-efficient fine-tuning (PEFT) methods have emerged as a
practical alternative. Instead of updating all parameters, PEFT methods
introduce a small number of trainable components while keeping the
pre-trained weights frozen. Among these, LoRA \cite{hu2021lora} has
become particularly popular due to its simplicity and strong empirical
performance across a wide range of tasks.

In this paper, we study PEFT methods in the context of ModernBERT
\cite{warner2024modernbert}, a recently released encoder-only transformer
that incorporates architectural improvements developed since the original
BERT release in 2018. ModernBERT supports longer input sequences (up to
8,192 tokens), uses rotary positional embeddings, and achieves state-of-the-art
performance on both text and code understanding benchmarks. Despite its
recent release, there is currently no published study on the behaviour
of PEFT methods applied to ModernBERT, motivating our investigation.

We address the following research questions:
\begin{enumerate}
    \item Can LoRA match full fine-tuning on ModernBERT with significantly
          fewer trainable parameters?
    \item How does LoRA rank affect accuracy and efficiency?
    \item Does LoRA performance degrade in low-data settings?
    \item How does IA\textsuperscript{3} compare to LoRA on ModernBERT
          across different tasks?
\end{enumerate}

% ─────────────────────────────────────────
\section{Background and Related Work}

\subsection{BERT and ModernBERT}

BERT \cite{devlin2019bert} introduced bidirectional pre-training of
transformers, achieving state-of-the-art results across a wide range of
NLP benchmarks. Subsequent work such as RoBERTa \cite{liu2019roberta}
improved BERT by optimizing pre-training hyperparameters and using more
data, while DeBERTa \cite{he2021deberta} introduced disentangled
attention mechanisms.

ModernBERT \cite{warner2024modernbert} represents a
significant modernization of the BERT architectures with an extended context window and a design optimized for downstream performance and efficiency on modern hardware. It was also trained on a substantially larger and more diverse corpus than the original BERT models. To support longer input sequences, ModernBERT interleaves local sliding-window attention with periodic global attention, extending the context window to 8,192 tokens. Compared to earlier BERT variants, ModernBERT additionally incorporates Rotary Position Embeddings (RoPE), GeGLU activations, bias-free linear layers, and native support for FlashAttention. 

ModernBERT achieves state-of-the-art or competitive performance on a wide range of downstream NLP benchmarks compared to its predecessors. However, because these architectural changes differentiate ModernBERT from earlier BERT variants, it remains unclear whether conclusions drawn from previous evaluations of Parameter-Efficient Fine-Tuning (PEFT) methods generalize to this architecture.

\subsection{Parameter-Efficient Fine-Tuning}

PEFT methods reduce the computational and memory cost of adapting large pre-trained models by keeping most of the model parameters frozen while updating only a small subset of trainable parameters \cite{hanParameterEfficientFineTuningLarge2024}. Depending on the method, these trainable parameters may be introduced as additional modules, low-rank weight updates, learned prompts, or scaling vectors integrated throughout the network. Despite optimizing only a small fraction of the model parameters, PEFT methods have been shown to achieve performance comparable to full fine-tuning on many downstream tasks.


LoRA \cite{hu2021lora} is a reparameterization-based PEFT method that introduces small trainable matrices into selected layers while keeping the original model weights frozen. It is based on the observation that weight updates during adaptation have a low intrinsic dimension, meaning that the changes required for fine-tuning can often be well approximated by a low-rank decomposition. Specifically, LoRA freezes the original parameters and inserts two low-rank matrices, $A$ and $B$, referred to as adapters, into existing weight matrices. The resulting learnable update is given by
$
\Delta W = \frac{\alpha}{r}BA,
$
where $B \in \mathbb{R}^{d \times r}$, $A \in \mathbb{R}^{r \times k}$, and $r \ll \min(d, k)$ is the bottleneck rank that controls the expressive capacity of the adaptation. The scaling factor $\frac{\alpha}{r}$ determines the effective magnitude of the update and helps maintain a consistent update scale across different choices of rank. By learning only these low-rank matrices instead of updating all model parameters, LoRA substantially reduces memory usage and computational cost, making fine-tuning significantly more efficient than full fine-tuning.

IA\textsuperscript{3} \cite{liu2022ia3} takes a different approach being a multiplicative PEFT method that adapts a pre-trained model by learning a small set of trainable scaling vectors while keeping all original model parameters frozen. Instead of introducing additional weight matrices, IA\textsuperscript{3} modulates internal activations via element-wise multiplicative rescaling. Concretely, learned vectors are injected into selected components of the network, including the key and value projections in the attention mechanism and the hidden activations in the feed-forward layers.

For a given hidden representation $h$, IA\textsuperscript{3} applies a learned scaling vector $l$ as
\[
h' = l \odot h,
\]
where $\odot$ denotes element-wise multiplication. In the attention module, this corresponds to
\[
k' = l_k \odot k, \quad v' = l_v \odot v,
\]
while in the feed-forward network, a separate vector $l_{\text{ffn}}$ rescales intermediate activations. These vectors are low-dimensional (one value per channel), making the number of trainable parameters extremely small compared to full fine-tuning or reparameterization-based methods such as LoRA. By selectively amplifying or suppressing feature dimensions through learned gating, IA\textsuperscript{3} achieves efficient task adaptation with minimal computational and memory overhead.

% ─────────────────────────────────────────
\section{Methodology}

\subsection{Base Model}

We use \texttt{answerdotai/ModernBERT-base} as the base model for all experiments. ModernBERT-base contains 149.6M
parameters. We freeze or update these base weights depending on the specific fine-tuning strategy, ensuring the starting representations remain identical across all trials.

\subsection{Datasets}

We evaluate on two GLUE \cite{wang2019glue} benchmarks:

\paragraph{SST-2} The Stanford Sentiment Treebank \cite{socher2013sst2}
is a binary sentiment classification task consisting of movie review
sentences labeled as positive or negative. We train on the standard split
(67,349 examples) and report accuracy on the validation set (872 examples).

\paragraph{QNLI} The Question Natural Language Inference dataset is derived from SQuAD. The model determines whether a given sentence entails the answer to an associated question. We use 104,743 training examples and report accuracy on the validation set of 5,463 examples.

\subsection{Fine-Tuning Methods}

We compare three fine-tuning strategies:

\paragraph{Full fine-tuning} All 149.6M model parameters are updated
during training. We use this as our upper-bound baseline.

\paragraph{LoRA} We apply Low-Rank Adaptation (LoRA) matrices into the query-key-value projections (\texttt{Wqkv}) of each transformer block, as well as the final sequence pooler projection layer. We evaluate ranks $r \in \{4, 8, 16\}$, setting $\alpha = 2r$ and applying a dropout rate of 0.1. For evaluating LoRA in low-data settings, the model is trained with 500, 1000, and 5000 training examples.

\paragraph{IA\textsuperscript{3}} We apply IA\textsuperscript{3} learned rescaling vectors to the query-key-value projections (\texttt{Wqkv}) of each transformer block. Furthermore, we apply feedforward rescaling vectors to the final sequence pooler projection layer (\texttt{dense}).

\subsection{Training Setup}

Models are trained using the HuggingFace \texttt{Trainer} API with the default AdamW optimiser. We maintain a constant learning rate of $2 \times 10^{-5}$ across all experiments, including both full fine-tuning and parameter-efficient methods. We use a training batch size of 16 and an evaluation batch size of 32. 

Model evaluation occurs at the end of each training epoch. We explicitly disable checkpoint saving; consequently, all reported validation accuracies reflect the model's state at the final training step rather than the peak validation performance. We record the total wall-clock training time to compare the computational overhead of each method. All experiments are executed on a single NVIDIA T4 GPU via Google Colab, with telemetry disabled and a fixed random seed of 42 to ensure reproducibility.

% ─────────────────────────────────────────
\section{Results}

\subsection{Full Fine-Tuning Baseline}

Table~\ref{tab:baseline} shows results for full fine-tuning across
different numbers of training epochs on SST-2.

\begin{table}[h]
\centering
\small
\begin{tabular}{lccc}
\toprule
\textbf{Epochs} & \textbf{Accuracy} & \textbf{Params} & \textbf{Time (min)} \\
\midrule
1 & 95.30 & 149.6M & 32.8 \\
2 & 94.95 & 149.6M & 72.7 \\
3 & 94.72 & 149.6M & 101.6 \\
\bottomrule
\end{tabular}
\caption{Full fine-tuning results on SST-2 validation set.}
\label{tab:baseline}
\end{table}

Notably, a single epoch achieves the highest accuracy (95.30\%), with
additional epochs slightly degrading performance, suggesting rapid
convergence and mild overfitting on subsequent epochs.

\subsection{LoRA Rank Ablation}

Table~\ref{tab:lora_rank} reports the effect of LoRA rank on SST-2.

\begin{table}[h]
\centering
\small
\begin{tabular}{lccc}
\toprule
\textbf{Rank} & \textbf{Accuracy} & \textbf{Params} & \textbf{Time (min)} \\
\midrule
$r=4$  & 92.20 & 278K  & 83.9 \\
$r=8$  & 92.43 & 554K  & 79.3 \\
$r=16$ & 93.00 & 1.1M  & 79.9 \\
\bottomrule
\end{tabular}
\caption{LoRA rank ablation on SST-2 (full training set).}
\label{tab:lora_rank}
\end{table}

LoRA with $r=16$ achieves 93.00\% accuracy, within 2.3 percentage points
of the full fine-tuning baseline, while using only 0.74\% of the
trainable parameters. The accuracy difference between $r=4$ and $r=16$
is small (0.80\%), suggesting rank is not the primary bottleneck for
this task.

\subsection{Low-Data Regime}

Table~\ref{tab:lowdata} shows the effect of training set size on LoRA
performance ($r=8$) on SST-2.

\begin{table}[h]
\centering
\small
\begin{tabular}{lcc}
\toprule
\textbf{Train Samples} & \textbf{Accuracy} & \textbf{Time (min)} \\
\midrule
500   & 47.71 & 1.1  \\
1,000 & 54.01 & 1.8  \\
5,000 & 60.09 & 6.4  \\
Full (67,349) & 93.00 & 78.7 \\
\bottomrule
\end{tabular}
\caption{LoRA ($r=8$) performance under varying training set sizes on SST-2.}
\label{tab:lowdata}
\end{table}

Performance degrades substantially with limited data, dropping to 60.09\%
with 5,000 samples --- barely above the majority class baseline for a
binary classification task. This suggests LoRA adapters on ModernBERT
require a sufficient volume of task-specific data to update meaningfully.

\subsection{IA\textsuperscript{3} vs LoRA Comparison}

% ── Fill in your results here once experiments are complete ──
% Example table structure below — replace ? with your values

\begin{table}[h]
\centering
\small
\setlength{\tabcolsep}{4pt}
\begin{tabular}{llccc}
\toprule
\textbf{Method} & \textbf{Dataset} & \textbf{Acc.} & \textbf{Params} & \textbf{Time} \\
\midrule
Full FT          & SST-2 & 95.30 & 149.6M & 32.8 \\
LoRA $r=8$       & SST-2 & 93.00 & 554K   & 79.3 \\
IA$^3$           & SST-2 & 82.34 & 53.0K  & 74.5 \\
\midrule
Full FT          & QNLI  & 50.54 & 149.6M & 55.7 \\
LoRA $r=8$       & QNLI  & 49.70 & 555K   & 133.3 \\
IA$^3$           & QNLI  & 81.19 & 53.8K  & 75.8 \\
\bottomrule
\end{tabular}
\caption{Comparison of fine-tuning methods across SST-2 and QNLI. Full FT = full fine-tuning.}
\label{tab:comparison}
\end{table}

The results on QNLI diverge sharply from those on SST-2. Both full
fine-tuning (50.54\%) and LoRA (49.70\%) fail to exceed chance-level
accuracy on this binary task, while IA$^3$ reaches 81.19\%, only 1.15
points below its own SST-2 result. We discuss this surprising reversal,
together with our suspected causes, in Section~\ref{sec:discussion}.

% ─────────────────────────────────────────
\section{Discussion}
\label{sec:discussion}

\subsection{LoRA is Effective but Dependent on Data}

On SST-2, LoRA behaves the way the PEFT literature would predict: it
adapts ModernBERT effectively for text classification, reaching within
2.3 points of full fine-tuning accuracy with a small fraction of
trainable parameters. This holds, however, only when sufficient training
data is available. The sharp accuracy drop in our low-data experiments
(47.71\% with 500 samples, barely above the majority class baseline)
suggests that LoRA's low-rank adapters need enough gradient signal across
many examples before they capture meaningful task specific structure.
With too few examples, the update direction is too noisy to settle on a
useful low-rank approximation, and the model essentially fails to move
away from chance.

\subsection{Rank Has Limited Impact}

On SST-2, the accuracy difference between $r=4$ and $r=16$ is small
(0.80 points), suggesting that for a comparatively easy task such as
binary sentiment classification, a low rank is already expressive enough
to capture what the model needs to learn. Practically, this means
practitioners working on similarly simple tasks may be able to use
smaller ranks and save memory without paying a real accuracy cost.
Whether this still holds for a harder task is precisely what our QNLI
experiments were meant to test. However, the answer obtained was rather unexpected.

\subsection{The QNLI Failure Interpretation}

The most unanticipated result, and the one most important to
discuss honestly, is that both full fine-tuning (50.54\%) and LoRA
(49.70\%) failed to learn QNLI at all: both sit essentially at the
50\% chance level of a binary classification task, while only IA$^3$
reached a respectable 81.19\%. This is not a case of one method
slightly underperforming another. In fact, two of our three methods, including
full fine-tuning, did not learn the task. This indicates that
the cause is unlikely to lie in LoRA, IA$^3$, or ModernBERT individually,
but rather in our own training setup.

This does not imply that LoRA is unsuitable for harder NLU tasks
in general. Most plausibly, the
problem traces back to a methodological choice that, in hindsight,
was suboptimal: we used the same fixed learning rate of $2 \times 10^{-5}$
for every method and every dataset, without any per task or per method
tuning, and without a learning rate warmup or schedule. QNLI is a larger
and structurally different task than SST-2: this means that it requires the model to
relate a question to a separate sentence rather than judge a single
sentence in isolation. Full fine-tuning of 149.6M parameters is
considerably more sensitive to the learning rate than a tiny
low-dimensional adapter, as the entire weight space is exposed to the
update rather than a constrained low-rank subspace \cite{hu2021lora}. A learning rate that is perfectly reasonable
for SST-2 may simply be unstable or insufficient for QNLI under full
fine-tuning, and the same instability likely propagated to LoRA, since
LoRA's adapters are inserted into the same attention projections that
full fine-tuning is updating directly. IA$^3$, by contrast, only learns
small element wise rescaling vectors with a much more constrained update
geometry \cite{liu2022ia3}, which may make it considerably more robust to a learning rate
that is not well tuned for the task.

A second, related possibility is that our training schedule itself was
too short or too abrupt for QNLI. We trained for a fixed three epochs
with no warmup and no learning rate decay, a setup carried over directly
from our SST-2 experiments without re-validation. SST-2 sentences are
short and the task is comparatively easy, so this schedule was
apparently sufficient there. QNLI pairs are longer and the task is
harder, and it is plausible that the model needed either more epochs, a
warmup phase to avoid early instability, or both. Unfortunately, we did not have the resources in our timeline to run a learning rate or schedule sweep on
QNLI for the scope of the project, which we consider a real gap in our
experimental design rather than a minor omission.

These results should be interpreted carefully: we cannot
claim, from our QNLI numbers alone, that LoRA underperforms full
fine-tuning or IA$^3$ on harder NLU tasks, because our full fine-tuning
baseline failed under the same conditions. What can be said is that
IA$^3$'s update mechanism was, in our specific setup, considerably more
robust to an untuned learning rate than either full fine-tuning or LoRA.
This is itself a useful, if accidental, finding, but it concerns
robustness to hyperparameter choice rather than which method is intrinsically stronger on harder tasks. A fair
comparison would require, at minimum, a small learning rate sweep
performed independently for each method on each dataset, and we list
this as the first priority for any continuation of this work.

\subsection{Training Time Considerations}

Training time in our experiments does not cleanly track parameter count,
and we think this is worth flagging rather than glossing over. On SST-2,
full fine-tuning for one epoch (32.8 minutes) was faster than LoRA for
three epochs (79.3 minutes), which on its own might suggest LoRA offers
no efficiency benefit. We believe this is largely an artifact of our
experimental design rather than a property of LoRA itself: we fixed the
number of epochs per method rather than the number of epochs per
comparison, so the full fine-tuning baseline that converged fastest was
also the one we happened to train for the fewest epochs. On QNLI the
picture is messier still. In fact, LoRA took 133.3 minutes, almost two and a
half times longer than full fine-tuning's 55.7 minutes on the same
data, which we cannot fully explain from our logs and suspect is
related to Colab's shared, variable GPU availability rather than a
genuine property of the method. We did not control for background load
on the T4 instances we used, and wall clock time on free-tier Colab is,
in our experience, noisy enough that small numbers of runs are not a
reliable basis for strong claims about relative speed. A more rigorous
comparison would fix the number of epochs across methods, repeat each
configuration multiple times, and report variance, which would be the basis for future research about this topic.

\subsection{Limitations}

We would rather list our limitations plainly than understate them.
First, our learning rate and training schedule were not tuned per
method or per dataset, and we believe this is the most likely cause of
the QNLI failures described above. This is the limitation we would fix
first given more time. Second, we evaluate on only two classification
tasks, both at the sentence level. Results may look completely different
on generation tasks or token level tasks such as named entity
recognition, which we did not have enough time to properly run. Third, we test only
ModernBERT-base, as the larger ModernBERT-large variant may behave
differently under the same PEFT methods. Fourth, we did not tune LoRA's
or IA$^3$'s own hyperparameters (LoRA's $\alpha$ and dropout, or which
modules each method targets) beyond the single configuration we report,
so the numbers in this paper should be read as a single operating point
rather than each method's best achievable performance. Furthermore, we
disabled checkpoint saving to keep runs lightweight on free-tier Colab,
which means our reported accuracies reflect the final training step
rather than the best validation checkpoint. For the QNLI full
fine-tuning and LoRA runs in particular, it is possible that an earlier
checkpoint performed better before training degraded, and we have no way
to check this retroactively. Finally, we do not compare against other
PEFT methods such as prefix tuning or AdaLoRA \cite{zhang2022adalora},
which would help establish whether the robustness we observed in IA$^3$
is specific to that method or shared by other multiplicative,
low-parameter approaches.

% ─────────────────────────────────────────
\section{Conclusion}

We present an empirical study of parameter-efficient fine-tuning applied
to ModernBERT, a state-of-the-art encoder-only transformer released in
December 2024. On SST-2, our results match expectations from the broader
PEFT literature: LoRA reaches within 2.3 points of full fine-tuning
accuracy while using 99.6\% fewer trainable parameters, rank has only a
small effect on accuracy, and performance degrades sharply once training
data becomes scarce.

Our QNLI experiments, however, did not go as planned, and the
honest reporting of that is as valuable as the SST-2 results
themselves. Both full fine-tuning and LoRA failed to learn the task
under our fixed, untuned learning rate and schedule, landing at
chance level accuracy, while IA$^3$ reached 81.19\% under the exact same
conditions. Rather than concluding that IA$^3$ is simply the better
method, we conclude that our comparison was not fair as designed: a
single learning rate chosen for SST-2 and reused without adjustment was
not adequate for a harder, structurally different task, and we could not correct this before reporting
results. We see real value in reporting this failure rather than
omitting the QNLI experiments or quietly re-running them with a better
learning rate after the fact, since it points to a concrete and easily
overlooked pitfall: comparing PEFT methods fairly requires tuning each
method's optimization setup independently per task, not reusing one
configuration that happened to work on an easier dataset.

Future work should, first and most directly, repeat the QNLI comparison
with a per method, per task learning rate sweep and a proper warmup
schedule to establish whether LoRA and full fine-tuning can in fact
learn the task under fairer conditions. Beyond that, extending this
study to token level and generation tasks, to ModernBERT-large, and to
additional PEFT methods such as AdaLoRA would give a much more complete
picture of when parameter efficient fine-tuning is and is not a safe
default choice for ModernBERT.

% ─────────────────────────────────────────
\bibliographystyle{acl_natbib}
\bibliography{references}

\end{document}
